\subsection{The genuine finite-boundary norm completion}\label{subsec:fc-line-affinoid-algebra}

\begin{definition}[The actual completed function ring]\label{def:fc-line-affinoid-algebra}
\lean{EnvelopeProperty.Berkovich.RelativeSpectrum.FCStandardAffinoid.rationalBoundaryFintype, EnvelopeProperty.Berkovich.RelativeSpectrum.FCStandardAffinoid.NormedRationalFunctions, EnvelopeProperty.Berkovich.RelativeSpectrum.FCStandardAffinoid.normedRationalFunctionsCommRing, EnvelopeProperty.Berkovich.RelativeSpectrum.FCStandardAffinoid.normedRationalEquiv, EnvelopeProperty.Berkovich.RelativeSpectrum.FCStandardAffinoid.rationalCoefficients, EnvelopeProperty.Berkovich.RelativeSpectrum.FCStandardAffinoid.normedRationalFunctionsAlgebra, EnvelopeProperty.Berkovich.RelativeSpectrum.FCStandardAffinoid.rationalBoundaryEvaluation, EnvelopeProperty.Berkovich.RelativeSpectrum.FCStandardAffinoid.normedRationalFunctionsSeminormedCommRing, EnvelopeProperty.Berkovich.RelativeSpectrum.FCStandardAffinoid.normedRationalFunctionsNormOneClass, EnvelopeProperty.Berkovich.RelativeSpectrum.FCStandardAffinoid.normedRationalFunctionsIsUltrametricDist, EnvelopeProperty.Berkovich.RelativeSpectrum.FCStandardAffinoid.nnnorm_normedRationalFunctions, EnvelopeProperty.Berkovich.RelativeSpectrum.FCStandardAffinoid.nnnorm_rationalCoefficients, EnvelopeProperty.Berkovich.RelativeSpectrum.FCStandardAffinoid.AnalyticFunctions, EnvelopeProperty.Berkovich.RelativeSpectrum.FCStandardAffinoid.analyticFunctionsNormOneClass, EnvelopeProperty.Berkovich.RelativeSpectrum.FCStandardAffinoid.analyticCoefficients, EnvelopeProperty.Berkovich.RelativeSpectrum.FCStandardAffinoid.analyticFunctionsAlgebra, EnvelopeProperty.Berkovich.RelativeSpectrum.FCStandardAffinoid.nnnorm_algebraMap_analyticFunctions, EnvelopeProperty.Berkovich.RelativeSpectrum.FCStandardAffinoid.rationalMap, EnvelopeProperty.Berkovich.RelativeSpectrum.FCStandardAffinoid.denseRange_rationalMap, EnvelopeProperty.Berkovich.RelativeSpectrum.FCStandardAffinoid.nnnorm_rationalMap}
\leanok
\uses{def:fc-line-rational-functions, thm:fc-rational-boundary-maximum}
Give the actual rational algebra $R_U$ the boundary seminorm
\[
 \|f\|_U=\max_{z\in\partial U}|f|_z.
\]
Its separated norm completion is the actual function ring $\mathcal O(U)$.
There is a dense original rational-function ring map $R_U\to\mathcal O(U)$,
and its norm is exactly the displayed boundary seminorm.
Its coefficient map is the original $F$-algebra coefficient map followed by
completion, and the norm of the image of $a\in F$ is exactly $|a|$.
\end{definition}
\begin{proof}
Evaluate $R_U$ in the finite product $\prod_{z\in\partial U}\mathcal H(z)$
of the original completed residue fields. Pull back its maximum norm along
this actual ring homomorphism. This gives a submultiplicative ultrametric
ring seminorm, with exactly the displayed value. The genuine separated
completion of this seminormed ring is a complete normed commutative ring;
the original rational map has dense range by the defining property of
completion and preserves its seminorm. Every original boundary point has its
prescribed value $|a|$ on the same coefficient $a$, so the finite-product norm
of that coefficient is $|a|$; completion retains that norm. The pulled-back
norm is ultrametric because the original completed residue fields are
ultrametric. This construction also covers the
zero-radius case: its rational seminorm can have a kernel, and the separated
completion identifies precisely elements at seminorm distance zero.
No faithfulness of that seminorm is assumed.
\end{proof}

\begin{lemma}[Actual completed evaluation]\label{lem:fc-line-completed-evaluation}
\lean{EnvelopeProperty.Berkovich.RelativeSpectrum.FCStandardAffinoid.nnnorm_rationalEvaluation_le, EnvelopeProperty.Berkovich.RelativeSpectrum.FCStandardAffinoid.normedRationalEvaluation, EnvelopeProperty.Berkovich.RelativeSpectrum.FCStandardAffinoid.lipschitz_normedRationalEvaluation, EnvelopeProperty.Berkovich.RelativeSpectrum.FCStandardAffinoid.evaluationToH, EnvelopeProperty.Berkovich.RelativeSpectrum.FCStandardAffinoid.evaluationToH_coe, EnvelopeProperty.Berkovich.RelativeSpectrum.FCStandardAffinoid.evaluationToH_rationalMap, EnvelopeProperty.Berkovich.RelativeSpectrum.FCStandardAffinoid.continuous_evaluationToH, EnvelopeProperty.Berkovich.RelativeSpectrum.FCStandardAffinoid.evaluationToH_unique, EnvelopeProperty.Berkovich.RelativeSpectrum.FCStandardAffinoid.nnnorm_evaluationToH_le, EnvelopeProperty.Berkovich.RelativeSpectrum.FCStandardAffinoid.analyticValue, EnvelopeProperty.Berkovich.RelativeSpectrum.FCStandardAffinoid.analyticValue_rationalMap, EnvelopeProperty.Berkovich.RelativeSpectrum.FCStandardAffinoid.analyticValue_le_nnnorm}
\leanok
\uses{def:fc-line-affinoid-algebra, thm:fc-rational-boundary-maximum}
For every original point $x\in U$, the original rational evaluation extends
uniquely to a continuous ring homomorphism
$\mathcal O(U)\to\mathcal H(x)$. It retains every original polynomial and
rational value and satisfies $|f|_x\le\|f\|_U$ for every completed function.
\end{lemma}
\begin{proof}
The actual rational boundary maximum theorem gives
$|f|_x\le\max_{z\in\partial U}|f|_z=\|f\|_U$ for every original rational
function. Apply this to differences to make the original completed-residue
field evaluator contractive. Its target is complete, so completion gives a
continuous ring extension with the original values. Density gives uniqueness.
The same bound on every completed element follows by continuity and density.
\end{proof}

\begin{theorem}[Exact finite-boundary norm and completed maximum modulus]\label{thm:fc-completed-boundary-maximum}
\lean{EnvelopeProperty.Berkovich.RelativeSpectrum.FCStandardAffinoid.nnnorm_eq_boundary_nnnorm, EnvelopeProperty.Berkovich.RelativeSpectrum.FCStandardAffinoid.nnnorm_eq_boundary_max, EnvelopeProperty.Berkovich.RelativeSpectrum.FCStandardAffinoid.analyticValue_exists_boundary_norm, EnvelopeProperty.Berkovich.RelativeSpectrum.FCStandardAffinoid.analyticValue_exists_boundary_max, EnvelopeProperty.Berkovich.RelativeSpectrum.FCStandardAffinoid.nnnorm_pow_analyticFunctions}
\leanok
\uses{def:fc-line-affinoid-algebra, lem:fc-line-completed-evaluation}
For every $f\in\mathcal O(U)$,
\[
 \|f\|_U=\max_{z\in\partial U}|f|_z.
\]
In particular, an actual Gauss boundary point realizes the norm, and every
original point of $U$ has value at most that same maximum.
Moreover $\|f^n\|_U=\|f\|_U^n$ for every nonnegative integer $n$.
\end{theorem}
\begin{proof}
For rational functions the equality is the definition of the pulled-back
finite-product norm. Both the completion norm and the norm of the finite
tuple of completed boundary evaluations are continuous functions on
$\mathcal O(U)$. They agree on the dense original rational functions, hence
everywhere. The boundary is a nonempty finite set, containing the original
outer Gauss point, so one original boundary point realizes this finite
maximum. Combine the equality with the preceding evaluation bound.
For power multiplicativity, submultiplicativity gives the upper bound.
Choose an original boundary point realizing $\|f\|_U$; evaluation there is
a ring homomorphism into a normed field, so it takes $f^n$ to a value of
norm $\|f\|_U^n$. The completed evaluation bound gives the reverse inequality.
\end{proof}
